\section{Justificación}

La transcripción automática de la música presenta un desafío al momento de procesar las señales digitales, de igual forma constituye un problema que toma gran relevancia al momento de recuperar información musical (MIR), esta tiene un gran potencial para transformar la manera en que se analiza, se interpreta y se accede al contenido musical. A pesar de los avances recientes en los ámbitos de machine learning, lo sistemas actuales no tienen la capacidad máxima de precisión y flexibilidad al comparar con un músico profesional, ya que estos tienen la capacidad de entender información polifónica mucho más amplia, estos escenarios polifónicos complejos que son escenarios donde convergen multiples voces, sonidos y narrativas sonoras.
